\section{课程背景与挑战}

UC Berkeley CS294-196 Fall 2024 的这讲由 Meta AI FAIR 的 Yuandong Tian 主讲，题为 ``Towards a Unified Framework of Neural and Symbolic Decision Making''，重点讨论如何把神经网络、搜索、规划与 agent 数据闭环统一到同一套决策框架里。
他以一条渐进趋势开场：随着任务难度升高，最好的模型性能曲线仍然单调下降，并没有出现弹性。
\textit{``This is actually in some sense a very worrisome kind of curve that even the best model has not solved this planning problem.''}
通过对多个 benchmark（智能规划、Soang 游戏、Maze 路径搜索）重复观察，同样的趋势一再出现：模型的分数被计算图里隐含的符号属性钳制住了，简单堆参数段时间内无法打开局面。
这为本讲的三条应对路径奠定了基调：单靠更大更贵的模型是不够的。

\subsection{规划能力中的曲线}

讲者在前 10 分钟内展示了一条针对复杂规划任务的评估曲线：虽然模型精度能改善，但收敛后的得分仍然偏低，而且所需的训练 token / 参数数量呈指数级上升。
这种现象反复在对话、自动规划、机器人导航里出现，问题的本质在于概率生成模型无法保证符号性质（约束、可行性）。
\begin{knowledgebox}{规划失败的症状}
LLM 在规划中最常见的症状包括：1）约束违规（如忽略边界或能耗限制）；2）长 horizon 时错误累积；3）没有对状态空间做精确的分叉管理。任何依赖精确性的问题如物流调度、合规审批都会被放大。
\end{knowledgebox}

\subsection{三条策略}

讲者提出三条修复路径：\textbf{Scaling}（加大 token/参数），\textbf{Hybrid}（LLM + Solver 混合），和更系统性的 \textbf{Compound AI}. 这不是互斥的选择，而是按照成本-思维跨度分层的架构，并且可以逐层叠加。
\begin{enumerate}
    \item \textbf{Scaling} 端通过更大的语言建模、更多训练数据和更多 compute 指望概率泛化，但对约束变量的结构性知识敏感度较低。
    \item \textbf{Hybrid} 端用 LLM 提供 candidate/action+heuristic，求解器负责 enforce constraints，减少了纯生成模型的不可控性。
    \item \textbf{Compound AI} 端把多个模块串联（search trace、multi-turn dialogue、constrained solver），形成可审计的反馈回路，方便分工与调试。
\end{enumerate}
\begin{figure}[H]
\centering
\includegraphics[width=\textwidth]{frame-solution-categories.jpg}
\caption{讲者列出三类解决策略，定位在不同的工程成本/思维跨度上。 \protect\footnotemark}
\end{figure}
\footnotetext{视频画面时间区间：00:09:10--00:09:18。}
\begin{importantbox}{三策略拆解}
Scaling 重仓资源，Hybrid 通过 LLM 提供启发式而让求解器负责搜索，Compound AI 则把多种模块串联起来形成内建反馈回路，既可调控质量也便于诊断。
\end{importantbox}
\begin{warningbox}{不要盲目相信 Scaling}
垂直依赖更大的模型意味着更长的训练周期、更难调试的过拟合、以及过高的推理功耗；如果看不到结构性改进，额外参数最终只会放大 ``看起来合理'' 的错误。
\end{warningbox}

\subsection{本章小结}

LLM 生成机制在规划问题上天然存在不确定性，需要借助符号结构来恢复可控性，三条策略提供了从工程到方法论上的不同切入点。

